\documentclass[11pt,a4paper]{article}
\usepackage[utf8]{inputenc}
\usepackage[T1]{fontenc}
\usepackage[portuguese]{babel}
\usepackage{geometry}
\geometry{margin=2.5cm}
\usepackage{listings}
\usepackage{xcolor}
\usepackage{amsmath}
\usepackage{amssymb}
\usepackage{hyperref}
\usepackage{enumitem}
\usepackage{tcolorbox}
\usepackage{tabularx}
\usepackage{booktabs}
\usepackage{graphicx}

\definecolor{codebg}{rgb}{0.95,0.95,0.95}
\definecolor{codegreen}{rgb}{0.1,0.5,0.1}
\definecolor{codepurple}{rgb}{0.5,0.1,0.5}
\definecolor{codegray}{rgb}{0.5,0.5,0.5}

\lstdefinestyle{python}{
    backgroundcolor=\color{codebg},
    commentstyle=\color{codegray}\itshape,
    keywordstyle=\color{codepurple}\bfseries,
    stringstyle=\color{codegreen},
    basicstyle=\ttfamily\small,
    breaklines=true,
    frame=single,
    language=Python,
    showstringspaces=false,
    tabsize=4,
    literate={á}{{\'a}}1 {ã}{{\~a}}1 {é}{{\'e}}1 {í}{{\'i}}1 {ó}{{\'o}}1 {õ}{{\~o}}1 {ú}{{\'u}}1 {ç}{{\c{c}}}1 {Á}{{\'A}}1 {Ã}{{\~A}}1 {É}{{\'E}}1 {Í}{{\'I}}1 {Ó}{{\'O}}1 {Õ}{{\~O}}1 {Ú}{{\'U}}1 {Ç}{{\c{C}}}1 {â}{{\^a}}1 {ê}{{\^e}}1 {ô}{{\^o}}1 {û}{{\^u}}1 {Â}{{\^A}}1 {Ê}{{\^E}}1 {Ô}{{\^O}}1 {Û}{{\^U}}1 {à}{{\`a}}1 {è}{{\`e}}1 {ì}{{\`i}}1 {ò}{{\`o}}1 {ù}{{\`u}}1 {ä}{{\"a}}1 {ö}{{\"o}}1 {Ä}{{\"A}}1 {Ö}{{\"O}}1
}
\lstset{style=python}

\newtcolorbox{objetivos}[1][]{
  colback=green!5!white,
  colframe=green!40!black,
  title=O que você vai aprender nesta página,
  fonttitle=\bfseries,
  #1
}

\newtcolorbox{exercicio}[1]{
  colback=blue!5!white,
  colframe=blue!40!black,
  title=#1,
  fonttitle=\bfseries
}

\newtcolorbox{solucao}{
  colback=yellow!5!white,
  colframe=orange!60!black,
  title=Solução proposta,
  fonttitle=\bfseries
}

\title{\textbf{Data Analysis with Python 2024--2025}\\Parte 7: Project Work\\\large Orientações de Solução e Comandos Úteis}
\author{Tradução e Adaptação: Universidade de Helsinque (MOOC.fi)}
\date{}

\begin{document}

\maketitle

\section*{Nota Importante}
Este material é uma tradução e adaptação baseada no curso \textit{Data Analysis with Python 2024-2025}, oferecido pela \textbf{Universidade de Helsinque (MOOC.fi)}. As orientações a seguir descrevem os padrões de execução e envio para os projetos finais.

\tableofcontents
\newpage

\section{Orientações de Solução - Capítulo 7}

Como o Capítulo 7 trata do \textbf{Trabalho de Projeto Final} (\emph{Project Work}), não há uma lista única de códigos fixos padrão, mas sim frentes práticas divididas nos três grandes temas escolhidos pelo estudante. Abaixo encontram-se os padrões de comandos e procedimentos recomendados para garantir o sucesso na entrega.

\subsection*{1. Validação Local de Exercícios}

Para testar o andamento das funções desenvolvidas dentro dos diretórios de projeto baixados pelo TMC, utilize o comando padrão no terminal:

\begin{solucao}
\begin{lstlisting}
# Executa os testes automatizados associados ao projeto ativo
tmc test
\end{lstlisting}
\end{solucao}

\subsection*{2. Como Executar Testes de Outros Alunos (Revisão por Pares)}

Para testar de forma segura o código de um colega sem correr o risco de corromper o seu próprio repositório de exercícios:

\begin{solucao}
\begin{lstlisting}
# 1. Mude para um diretorio temporario no sistema Unix
cd /tmp

# 2. Baixe os arquivos de teste originais da plataforma
tmc download -a hy-data-analysis-with-python-2020

# 3. Substitua o notebook do estudante na pasta correspondente
# Exemplo para Regressao:
# part08-e01_regression/src/project_notebook_regression_analysis.ipynb

# 4. Execute a bateria de testes isoladamente
tmc test part08-e01_regression
\end{lstlisting}
\end{solucao}

\subsection*{3. Envio de Tarefas via Terminal (TMC CLI)}

Caso prefira utilizar a interface de linha de comando para gerar o link de compartilhamento e submissão por \emph{paste}:

\begin{solucao}
\begin{lstlisting}
# Entre na pasta do exercicio do projeto e execute:
tmc paste
\end{lstlisting}
\end{solucao}

O terminal retornará um link URL oficial (iniciado com \url{https://tmc.mooc.fi/paste/}) que deverá ser copiado e submetido no formulário de acompanhamento da universidade.

\end{document}
